\documentclass[times,final,authoryear]{elsarticle}

\usepackage{medima}
\usepackage{framed,multirow}
\usepackage{graphicx}
\usepackage{booktabs}
\usepackage{subcaption}
\usepackage{tabularx}
\usepackage{threeparttable}

\usepackage{amssymb}
\usepackage{amsmath}
\usepackage{latexsym}
\usepackage[left]{lineno}

\usepackage{tikz}
\usetikzlibrary{
    calc,
    positioning,
    arrows.meta,
    shapes.geometric,
    shapes.symbols
}
\usepackage{url}
\usepackage{xcolor}

\usepackage{hyperref}
\hypersetup{hidelinks}

\definecolor{newcolor}{rgb}{.8,.349,.1}

\journal{Agricultural Systems}

\makeatletter
\gdef\@jnlurl{https://www.sciencedirect.com/journal/agricultural-systems}
\def\@issn{Agricultural Systems}

\def\elsjnllogoboxb#1#2{\minipage[b][#2][b]{#1}\mbox{}\endminipage}
\makeatother

\begin{document}

\verso{Martins \textit{et~al.}}

\begin{frontmatter}

\title{Assessing Context-Conditioned Soybean Sowing Date and Supplemental Irrigation Decisions Using DSSAT and Reinforcement Learning}

\author[2,6,7,8]{Marcella  Scoczynski Ribeiro \snm{Martins}} %0000-0002-5716-4968
\author[1]{Gabrielly de Queiroz \snm{Pereira}} %orcid: https://orcid.org/0000-0003-2663-0783 gqpereira@uepg.br
\author[2]{Hilkija Gaïus \snm{Tosso}} %orcid https://orcid.org/0000-0003-1109-6813 hgaius@hotmail.com
\author[1,8]{Maria Salete Marcon Gomes \snm{Vaz}} %orcid https://orcid.org/0000-0001-9172-1863 salete@uepg.br
\author[1,8]{José Carlos Ferreira da \snm{Rocha}} %orcid https://orcid.org/0000-0002-4050-281X jrocha@uepg.br
\author[4,5,6]{Viviane M. \snm{Gomes Pacheco}} %0000-0003-3693-4377 viviane.gomes@ifg.edu.br
\author[4]{Cloves Gonçalves \snm{Rodrigues}} %0000-0003-0140-9847 cloves@pucgoias.edu.br
\author[3]{Antonio~Paulo~\snm{Coimbra}} %0000-0003-3780-2405 acoimbra@isr.uc.pt
\author[3,5,6]{Wesley Pacheco \snm{Calixto}\corref{cor1}} %0000-0002-1928-4432 wesley.pacheco@isr.uc.pt
\cortext[cor1]{Corresponding author.}
\ead{wesley.pacheco@isr.uc.pt}

\address[1]{State University of Ponta Grossa, Ponta Grossa, Parana/Brazil}
\address[2]{Federal University of Technology - Parana, Brazil}
\address[3]{Institute of Systems and Robotics, Coimbra University, Coimbra/Portugal}
\address[4]{Polytechnic and Arts School, Pontifical Catholic University of Goias, Goiania, Goias/Brazil PC~74.605-010}
\address[5]{Electrical, Mechanical \& Computer Engineering School, Federal University of Goias, Goiania, Goias/Brazil PC~74.605-010}
\address[6]{Experimental \& Technological Research and Study Group, Federal Institute of Goias, Goiania, Goias,  Brazil}
\address[7]{Federal University of Technology - Parana, Graduate Program in Electrical Engineering, Ponta Grossa, Brazil}
\address[8]{State University of Ponta Grossa, Postgraduate Program in Applied Computing, Ponta Grossa, Parana/Brazil}




\begin{abstract}
\textbf{CONTEXT:} Sowing date and supplemental irrigation are coupled decisions in predominantly rainfed soybean systems and are chosen before seasonal weather is known. A context-conditioned rule must be evaluated against comparators that receive the same information and simulation budget, on seasons not used to build it.

\textbf{OBJECTIVE:} To formulate soybean sowing date and irrigation-rule selection as a one-step contextual decision problem coupled to DSSAT-CROPGRO-Soybean, and to test whether a policy trained with Proximal Policy Optimization (PPO) on information available on September 1 adds value over an optimized fixed rule and a simple contextual rule in Castro, Paraná, Brazil, under two weather records of different length and quality.

\textbf{METHODS:} PPO selected sowing date and three parameters of a causal irrigation rule from six pre-season variables (rainfall, temperature and radiation before the decision date, and the Oceanic Ni\~no Index). Each action was evaluated by a full DSSAT run with raw simulated yield after a soil-water spin-up. Experiment 1 used the INMET A819 station (2006--2024) with a chronological split and five seeds. Experiment 2 used the NASA POWER daily series for the same coordinates (1984--2024) in a registered rolling-origin design with five folds, three seeds per fold and 20 test seasons. All rules were compared season by season on a common grid of 602 candidate rules, together with a retrospective oracle.

\textbf{RESULTS:} With the station record, PPO exceeded the optimized fixed rule on training and validation seasons but produced 843 kg ha$^{-1}$ less on the 2021--2024 test block, because missing station temperature placed two test contexts outside the training range and the four validation seasons required no irrigation. With the 41-season NASA POWER record, PPO reached a mean test yield of 5,590 kg ha$^{-1}$ with 27.4 mm of irrigation. The paired difference with the optimized fixed rule was $-92$ kg ha$^{-1}$ (95\% season-bootstrap interval $-226$ to 21), and that with the rainfed October 15 calendar was 366 kg ha$^{-1}$ (106 to 672). PPO irrigated 75 mm on average in the six test seasons with a yield response to irrigation and 7 mm in the other fourteen.

\textbf{CONCLUSIONS:} With a longer and homogeneous weather record, PPO learned a context-conditioned policy that matched the optimized fixed rule within the season-to-season uncertainty and exceeded the calendar and constructed schedules, but did not exceed the fixed rule, whose in-season trigger already adapts irrigation to the realized season. The length and quality of the weather record determined out-of-sample performance more than the learning algorithm.

\end{abstract}


\begin{keyword}
Soybean production \sep Reinforcement learning\sep DSSAT Simulations \sep Proximal Policy Optimization \sep Decision analysis \sep Water allocation
\end{keyword}

\end{frontmatter}

\linenumbers

%===============================================================
\section{Introduction}
\label{introduction}

In soybean production in southern Brazil, water deficit is the main cause of the yield gap and irrigation is among the remedies proposed to close it \citep{sentelhas2015yieldgap}. A farmer chooses, once per season, a sowing date inside a window of roughly three months and, where equipment exists, a rule for supplemental irrigation. The outcome of that choice depends on the interaction between seasonal weather, soil water-holding capacity, cultivar phenology, and the amount of water available for irrigation, none of which is fully known at sowing. A fixed sowing calendar ignores this interaction. A decision rule that conditions sowing date and water use on information available before the season could, in principle, exploit it. Whether such a rule delivers a measurable benefit over a well-chosen calendar, and under which water restrictions within the Castro production context, is a question of system behaviour that agricultural systems modelling was developed to address \citep{jones2017history}. It is also a question in which the value of pre-season climate information is itself part of the analysis \citep{hansen2002climate}.

Agricultural management requires selecting actions whose consequences depend on interacting environmental and biological conditions. Sowing date, irrigation, cultivar traits, soil properties, and weather jointly affect crop development and final yield; for this reason, management design is better treated as a decision and optimization problem than as yield prediction under a fixed strategy \citep{tao2022optimizing, Xiao2024}. Recent work in agricultural systems has also emphasized that digital decision systems are not defined only by predictive accuracy: the objective function, its weights, constraints, and trade-off rules shape what the system is designed to optimize \citep{carolan2026objective}. This point is directly relevant to irrigation and sowing decisions, because increasing yield, reducing water use, limiting risk, and maintaining operational feasibility can lead to different preferred actions. In soybean production, sowing date is particularly important because it determines the environmental conditions experienced during critical phenological stages and can also affect the feasibility and productivity of subsequent crops within the same agricultural year \citep{debangshi2025plantingdates, bigolin2025latesoybean, zhang2021modissoybean}.

Process-based crop simulation models provide a suitable computational basis for this type of analysis. Unlike purely empirical approaches that estimate yield directly from historical observations, these models represent physiological and environmental processes governing crop growth and development. The Decision Support System for Agrotechnology Transfer (DSSAT) is one of the most widely used crop-modeling platforms and enables the simulation of crop responses to weather, soil, genotype, and management conditions \citep{ovando2018dssatsoybean, singh2024dssatsoybean}. Previous soybean applications have examined responses to temperature, precipitation, sowing date, and weather-data sources \citep{battisti2017soybeanmodels, battisti2018climatesensitivity, ovando2018dssatsoybean}, typically for a predefined set of treatments. As the number of decision variables grows, exhaustive evaluation becomes less efficient. Applications of an established crop model to a decision problem also presuppose that the model has been calibrated and tested for the case in question. When soil and cultivar parameters are generic, the simulated response carries an unquantified systematic error that any downstream optimization inherits.

Optimization has become a common framework for agricultural water-resource management because it allows management scenarios to be compared under constraints and competing objectives. Recent reviews of agricultural water and reservoir management describe the use of simulation--optimization models, multiobjective optimization, and machine-learning surrogate models to reduce computational cost when evaluating management alternatives \citep{kumari2026optimization}. In parallel, Agriculture 5.0 and drone-as-a-service systems are expanding the availability of data streams for crop monitoring, soil assessment, and automated farm operations \citep{adel2026drones}. These data streams can improve the context available to decision models, but they do not by themselves answer whether a context-aware policy outperforms a well-selected fixed rule under the same information and simulation budget.

Machine learning has been used to forecast soybean yield \citep{vonbloh2023mlsoybean, schwalbert2020soybeanforecast}. Forecasting, however, estimates the outcome expected under a given set of conditions or management assumptions, while management optimization requires choosing among alternative actions. Reinforcement learning addresses this second problem by learning a policy through repeated interaction with an environment, which may be a process-based crop simulator. In irrigation scheduling, \citet{kelly2024assessing} showed with AquaCrop-OSPy that deep RL provided its clearest advantage under strong water limitation, while optimized heuristic policies were difficult to outperform under wetter or more variable rainfall conditions. Using gym-DSSAT, \citet{balderas2025comparative} evaluated reinforcement learning algorithms for crop management tasks involving fertilization and irrigation. For planting-date adaptation, \citet{noa2026hybrid} combined a process-based crop model, a machine-learning emulator, and reinforcement learning to explore adaptive corn planting policies. The present study differs from these works by focusing on a joint pre-season sowing-date and irrigation-rule decision for soybean in Castro, with DSSAT executed directly for the evaluated management scenarios.

When the whole management strategy is fixed before the season, the decision problem has a single step and is a contextual optimization problem rather than a sequential control problem. Policy-gradient methods such as PPO remain applicable because they learn a stochastic mapping from seasonal context to a continuous action vector. Their use, however, must be judged against an optimized fixed rule, against context-free global optimizers, and against a simple contextual rule that receives the same information under the same simulation budget. In this study, PPO is one mechanism of context-conditioned policy search. The object of the study is the value of the context and of the management objective, not the algorithm alone.


Among the closest studies identified \citep{kelly2024assessing, balderas2025comparative, noa2026hybrid}, and in light of recent reviews and conceptual work in agricultural systems \citep{carolan2026objective, kumari2026optimization, adel2026drones}, we found no evaluation of a joint pre-season sowing-and-irrigation decision for soybean in which the comparators receive the same information and the same simulation budget as the learned policy, in which a simple contextual rule is included, and in which out-of-sample performance is examined under weather records of different length and quality. The novelty claimed here is restricted to these verifiable axes within the Castro case study. The central question is: \emph{under which conditions of climate, soil water storage, and water restriction does a contextual sowing-and-irrigation policy add value relative to a well-optimized fixed rule in Castro?} The question admits a null or negative answer, and such an answer is informative for the design of decision systems when its causes can be identified.

Three directional expectations were defined before the final analysis. First, the value of adaptation may interact with water availability and with the restriction on irrigation water, without a presumed monotonic direction. Second, the storage capacity of the soil may modulate the value of adaptation in a season-dependent way. Third, the value of pre-season information is related to its observed predictive skill for the coming season, measured before the policy is evaluated, without imposing an a priori ordering among climatology, pre-season observations, and seasonal climate indices. Each expectation has the unit of comparison of Section~\ref{subbaselines} and the pre-specified criterion of Section~\ref{subdesign}; the Conclusion reports the estimated direction and magnitude rather than confirmed or refuted hypotheses.

The study addresses five analytical objectives: i) to quantify the performance of a context-conditioned sowing-and-irrigation policy relative to an optimized fixed rule in Castro seasons not used for model development; ii) to determine whether the value of contextual adaptation persists when the policy is compared with a simple contextual rule receiving the same pre-season information and with constructed reference schedules; iii) to quantify how contextual adaptation modifies the productivity--water relation and the yield of the least favourable seasons; iv) to identify the conditions under which pre-season context changes the selected management action, including the predictive skill of each context variable; and v) to determine how the length and quality of the weather record used to train and evaluate the policy affect these results.

The intended contribution is to quantify, for Castro, whether context-conditioned soybean sowing and irrigation-rule selection provides value beyond a well-selected fixed rule, and to identify which properties of the data and of the evaluation design determine the answer. The contribution is methodological and diagnostic. It documents a reproducible DSSAT--policy coupling with raw simulated yield, a comparator structure in which every rule is evaluated on the same simulated seasons, two experiments that differ in the weather source (a station record of 19 seasons with gaps and a gridded record of 41 seasons without gaps), a registered rolling-origin evaluation over 20 test seasons, and learning curves and checkpoint selection for every training run. Spatial transfer to other municipalities is outside the scope of this study.

% =====================================================================
\section{Materials and methods}
\label{metodologia}

%%============================================================
\subsection{System boundary and study environment}
\label{subsystem}

The system studied is a predominantly rainfed soybean field with optional supplemental irrigation, represented for Castro (Paraná, Brazil) by one weather series, one soil profile, and one cultivar. Its spatial boundary is the field (one soil, one management); its temporal boundary is the growing season from sowing to physiological maturity, preceded by a soil-water spin-up that starts on August 1 of the sowing year (Section~\ref{subcropmodel}). The components and their interactions are: (i) the seasonal weather (rainfall, temperature, radiation), uncertain at sowing; (ii) the soil water balance, which converts rainfall and irrigation into water available to the crop and depends on the profile and on the antecedent moisture; (iii) the crop, whose phenology and yield formation respond to the timing of water and thermal stress relative to the sowing date; (iv) the management decision, taken once before the season, consisting of a sowing date within a window and a rule for supplemental irrigation; and (v) the water resource, bounded by a seasonal cap. The unit of decision is one season of one field. Nitrogen is assumed non-limiting (Section~\ref{subcropmodel}), so the conclusions do not extend to nutrient-constrained systems.

\emph{Decision instant and information horizons}: the operational decision date is fixed at $t_0=$ September 1 of the sowing year. The policy receives the context and emits the action before the sowing window opens on September 15. The admissible weather information ends on August 31, and the admissible climate-index information is the last value published before $t_0$ (Section~\ref{subdecision}). The decision is not revised after $t_0$; the irrigation rule then uses daily weather during the season only through its prescribed causal schedule, and DSSAT simulates from the spin-up date to physiological maturity. Information reaches the decision maker at two moments: before $t_0$ (admissible in the context) and during the season (used by the simulator, by the causal irrigation rule, and, in a labelled retrospective diagnostic, by the oracle).

The site coordinates are those of the INMET automatic station A819 (latitude $-24.78$, longitude $-50.05$, altitude 1,016 m). The crop is represented by the generic maturity-group-5 cultivar \texttt{990005} of \texttt{SBGRO048.CUL}, with plant population 30 m$^{-2}$, row spacing 45 cm, and sowing depth 5 cm, on a synthetic clayey profile (Table~\ref{tab_environment_a_current}). No field phenology, field yield, or soil-water observations are available for this environment, and the crop-model outputs are therefore treated as a simulation experiment and not as a field prediction. No claim is made for spatial generalization beyond Castro.

\begin{table}[!h]
\centering
\scriptsize
\caption{Executable environment of the two experiments. Field validation metrics are not estimated because no field phenology, yield, or soil-water observations are available for this environment.}
\label{tab_environment_a_current}
\begin{tabular}{p{2.3cm}p{2.6cm}p{1.6cm}p{1.0cm}p{1.4cm}p{1.1cm}p{3.6cm}}
\hline
Experiment & Weather source & Soil & AWC (mm) & Cultivar & Field obs. & Seasons and split \\
\hline
1 (station) & INMET A819, hourly, quality-controlled & synthetic clayey & 198 & 990005, MG 5 & 0 & 2006--2024; train 2006--2016, valid 2017--2020, test 2021--2024 \\
2 (gridded) & NASA POWER daily, MERRA-2 and CERES & synthetic clayey & 198 & 990005, MG 5 & 0 & 1984--2024; rolling origin, test 2005--2024 (Table~\ref{tab_folds}) \\
\hline
\end{tabular}
\end{table}

DSSAT-CSM with the CROPGRO-Soybean module is the process representation. For each candidate strategy a complete treatment is generated (daily weather, soil profile, initial soil water, cultivar, sowing date, population and spacing, prescribed irrigation events, simulation controls), the season is simulated, and grain yield and seasonal irrigation are returned: $\mathbf{z} = f_{\text{DSSAT}}(y,a,\psi)$, where $y$ is the season, $a$ the management action, and $\psi$ the fixed soil, cultivar, and model parameters. Simulations are executed through \texttt{DSSATTools}; the policy search uses Stable-Baselines3 and PyTorch (versions in Section~\ref{subrepro}).

Three concepts are kept distinct throughout: \emph{crop-model calibration}, the adjustment of DSSAT soil and cultivar parameters against field-scale observations, which was not possible with the available data; \emph{statistical correction}, an external transformation between simulated output and the municipal series, used in a previous diagnostic configuration and removed from the objective (Section~\ref{subobjective}); and \emph{policy performance}, the primary outcome, yield and water use of a management rule relative to baselines in seasons not used for development.

%%============================================================
\subsection{DSSAT configuration and crop-model validation status}
\label{subcropmodel}

The executable workflow does not contain field data for crop-model calibration or independent validation. No measured Castro soil profile, cultivar-specific phenology, field or trial yield, or observed soil-water series is available. Therefore, this manuscript does not report DSSAT calibration metrics such as phenology RMSE, yield RMSE, NRMSE, bias, concordance coefficient, or soil-water error. The DSSAT outputs are simulations of a generic cultivar on a synthetic profile under the weather of each series; they define the response surface on which the management rules are compared, not an agronomically validated prediction of field yield.

The soil is a synthetic clayey profile created in \texttt{src/dssat\_adapter.py}, with five layers to 150 cm and 198 mm of available water between lower limit and drained upper limit; its layer parameters are exported to Supplementary Table S2. Simulation controls are: soil-water balance on (\texttt{WATER = Y}), nitrogen dynamics off (\texttt{NITRO = N}), no fertilization, sowing prescribed (\texttt{PLANT = R}), and irrigation prescribed only when events exist (\texttt{IRRIG = R}; otherwise \texttt{IRRIG = N}). Nitrogen is therefore assumed non-limiting by configuration.

\emph{Soil-water spin-up and simulation window.} Each simulation starts on August 1 of the sowing year with every soil layer at its drained upper limit and runs through the sowing date to physiological maturity. The weather file passed to DSSAT covers August 1 of the sowing year to June 30 of the following year, so that late sowings have weather until maturity. The soil water at sowing is therefore produced by the August--sowing weather rather than set to a default value on the sowing day.

The sowing action window is September 15--December 20. The script \texttt{src/planting\_window\_audit.py} records the MAPA soybean ZARC ordinance for Paraná, safra 2024/2025 (Portaria SPA/MAPA no. 118; \url{https://www.gov.br/agricultura/pt-br/assuntos/riscos-seguro/programa-nacional-de-zoneamento-agricola-de-risco-climatico/portarias/safra-vigente/parana/PORTN118SOJAPR.ret2.pdf}), extracts the general cycle classes reported in the ordinance (100, 115 and 130 days) and the AD1--AD6 available-water classes, and checks the sowing dates against the coded window. The PDF text does not contain the municipality-specific Castro table; the ordinance directs users to the ZARC Oficial panel for municipality-, soil- and group-specific recommendations.

%%============================================================
\subsection{Data acquisition}
\label{subdata}

\subsubsection{Station weather: INMET A819}
\label{subweather}

Hourly meteorological observations of INMET automatic station A819 (Castro) were downloaded from the INMET historical data portal (\url{https://portal.inmet.gov.br/dadoshistoricos}, accessed on 26 September 2026) as annual CSV files from July 2006 to May 2026; the 2022 records and the first day of the series were taken from the complementary station export \path{dados_A819_H_2006-07-08_2026-01-01.csv}, used only to fill absent timestamps. The consolidated hourly file contains 174,432 records from 8 July 2006 00:00 to 31 May 2026 23:00. Sentinel values are converted to missing; values $\leq -90$, negative precipitation, temperatures outside $-20$ to $50\,^{\circ}$C, and radiation below zero or above 6000 kJ m$^{-2}$ h$^{-1}$ are removed. Daily precipitation and radiation are hourly sums (radiation converted to MJ m$^{-2}$); daily maximum and minimum temperature are hourly extremes. Days with fewer than 18 valid precipitation or temperature hours, or fewer than 6 valid radiation hours, are set to missing. Missing precipitation is filled with the day-of-year climatological mean (monthly mean as fallback); missing or sub-0.5 MJ m$^{-2}$ d$^{-1}$ radiation is estimated by the Hargreaves method and bounded to 1--35 MJ m$^{-2}$ d$^{-1}$; missing extremes are replaced by the daily mean and interpolated. The weather-completeness table by season is exported to \texttt{weather\_completeness\_by\_season.csv}, and the sensitivity of seasonal rainfall to the imputation strategy is computed by \texttt{src/weather\_imputation\_sensitivity.py}.

\subsubsection{Gridded weather: NASA POWER}
\label{subpower}

Daily weather for the station coordinates was obtained from the NASA Prediction Of Worldwide Energy Resources (POWER) project through its daily point API (\url{https://power.larc.nasa.gov/api/temporal/daily/point}), with the request parameters \texttt{community=AG}, \texttt{latitude=-24.78}, \texttt{longitude=-50.0464}, \texttt{start=19810101}, \texttt{end=20260531}, \texttt{format=CSV} and \texttt{time-standard=LST}, and the variables \texttt{PRECTOTCORR} (MERRA-2 bias-corrected precipitation, mm d$^{-1}$), \texttt{T2M}, \texttt{T2M\_MAX} and \texttt{T2M\_MIN} (MERRA-2 air temperature at 2 m, $^{\circ}$C), and \texttt{ALLSKY\_SFC\_SW\_DWN} (CERES/GEWEX all-sky surface shortwave irradiance, MJ m$^{-2}$ d$^{-1}$). The request was made on 7 October 2026 and returned 16,587 days. The file is stored without modification as \path{data/external/nasa_power/power_castro_daily_1981_2026.csv}. Values equal to $-999$ denote missing data; they occur only for radiation in 1981--1983, so the seasons used start in 1984. The meteorological variables refer to a $0.5^{\circ} \times 0.625^{\circ}$ grid cell with a mean elevation of 926 m; precipitation is a reanalysis product corrected with gauge-based data and is not an observation at the field. No imputation was applied to this series.

\emph{Agreement with the station.} The script \texttt{src/power\_station\_agreement.py} compares the two series on the days on which the station passes the daily validity rule of Section~\ref{subweather}. For each variable it reports the number of days, the means, the mean difference (station minus POWER), the mean absolute error, the root mean square error, and the Pearson correlation; it also compares September--April rainfall totals for the seasons with at least 90\% valid station rainfall days and at least 240 days in the window.

\subsubsection{Climate index}
\label{suboni}

The Oceanic Ni\~no Index (ONI; three-month running mean of ERSST sea-surface temperature anomalies in the Ni\~no 3.4 region) was obtained from NOAA/PSL (\url{https://psl.noaa.gov/data/correlation/oni.data}, file version of 4 August 2026 based on ERSST v6) and stored as \path{data/external/oni_psl_noaa.txt}. The May--July value is published in early August and is available on $t_0$; the June--August value is published after $t_0$ and is not used.

\subsubsection{Municipal yield series (contextual verification only)}
\label{subsidra}

Observed soybean yields are obtained from Table 1612 of the IBGE Automatic Recovery System (SIDRA), available at \url{https://sidra.ibge.gov.br/tabela/1612} (accessed on 26 September 2026); only the municipal average soybean yield of Castro (kg ha$^{-1}$) is used. The series carries a technological trend unrelated to the weather of each season, and its unit of observation (municipal annual average over many fields, cultivars, and management levels) does not coincide with the unit of simulation (one field, one management). It is used as contextual information on the range of yields and enters neither DSSAT nor the policy objective. A season is identified by its sowing year $y$ (sowing September--December of $y$, harvest January--April of $y+1$). Following the IBGE PAM rule that assigns production to the civil year in which most production is registered (\url{https://www.ibge.gov.br/biblioteca/visualizacao/periodicos/66/pam_2023_v50_br_notas_tecnicas.pdf}), season $y$ corresponds to the SIDRA value of year $y+1$; the audit table \texttt{season\_sidra\_pairing\_audit.csv} stores both candidate pairings.

%%============================================================
\subsection{Decision formulation and information set}
\label{subdecision}

The decision is treated as a one-step contextual problem: on the decision date $t_0$ the policy observes a context $s$, chooses an action $a$, the season is simulated, a terminal outcome $r(s,a)$ is returned, and the episode ends. No temporal credit assignment is involved. The term \emph{context-conditioned policy search} is used for the method; reinforcement learning refers to the PPO implementation used to search the stochastic policy.

\emph{Information set (context).} The context contains six variables available on $t_0$:
\begin{equation}
s = \big(P_{30},\; P_{90},\; \bar{T}_{30},\; \bar{R}_{90},\; \mathrm{ONI}_{\mathrm{MJJ}},\; \mathrm{ONI}_{\mathrm{MJJ}}-\mathrm{ONI}_{\mathrm{FMA}}\big),
\label{eq_context}
\end{equation}
where $P_{30}$ and $P_{90}$ are rainfall accumulated over the 30 and 90 days ending on August 31 (mm), $\bar{T}_{30}$ is the mean air temperature of the 30 days ending on August 31 ($^{\circ}$C), $\bar{R}_{90}$ is the mean daily global radiation of the 90 days ending on August 31 (MJ m$^{-2}$ d$^{-1}$), and $\mathrm{ONI}_{\mathrm{MJJ}}$ is the ONI for May--July of the sowing year, with its change since February--April as an indicator of the ENSO tendency. The weather variables are computed from the series used in each experiment. Each variable is standardized with the mean and standard deviation of the training seasons of the run; the scaler is stored with each trained model and applied unchanged to validation and test contexts. Realized weather after $t_0$ is never part of the context.

\emph{Action.} The policy outputs four continuous values in $[-1,1]$ mapped linearly to: sowing offset $d \in \{0,\dots,96\}$ days after September 15 (rounded to the nearest day); irrigation trigger $\theta \in [10, 150]$ mm (rounded to 1 mm); depth per event $l \in [0, 35]$ mm (rounded to 0.5 mm); and seasonal cap $L \in [20, 260]$ mm (rounded to 5 mm). Events with depth below 5 mm are not applied, so that rainfed management corresponds to a region of the continuous action space rather than to a single corner; when no event is scheduled the environment passes an empty schedule to DSSAT and $W = 0$. Rounding makes repeated actions map to identical DSSAT treatments, which are then read from an in-memory cache.

\emph{Irrigation rule.} The three irrigation parameters define a parameterized schedule computed from a proxy water balance on the daily weather from the sowing date, $B_{d+1} = \max(0, B_d + E_d - P_d)$ with $E_d = 0.16\,R_d + 0.08\,\max(\bar{T}_d - 10, 0)$ (mm d$^{-1}$, $R_d$ in MJ m$^{-2}$ d$^{-1}$, $\bar{T}_d$ in $^{\circ}$C), inspected every 4 days, the return interval of a centre pivot under supplemental irrigation: if $B_d \geq \theta$ and the accumulated irrigation is below $L$, an event of $\min(l, L - \sum W)$ mm is scheduled on day $d$ and $B_d$ is reduced by that depth. The decision to irrigate at the start of day $d$ uses $B_d$, computed exclusively from the weather observed up to the end of day $d-1$. The rule is evaluated over the 150 days after sowing, the events are passed to DSSAT as prescribed irrigation, and the rule is not a feedback on the simulated soil water. The proxy coefficients were not calibrated against reference evapotranspiration or observed soil-water deficit. The upper bound of $\theta$ (150 mm, about three quarters of the available water of the profile) was set on the development seasons of Experiment 1, where triggers $\leq 90$ mm irrigated in seasons without a yield response.

%%============================================================
\subsection{Policies and baselines}
\label{subbaselines}

Comparators are separated by the information they use, so that the value of a good global rule, the value of context, and the value of the algorithm can be attributed. All comparators except PPO are evaluated on a common grid of 602 candidate rules: 13 sowing offsets (0, 8, \dots, 88, 96 days after September 15) combined with rainfed management or with 45 irrigation rules ($\theta \in \{30, 60, 90, 120, 150\}$ mm, $l \in \{15, 25, 35\}$ mm, $L \in \{80, 160, 240\}$ mm), plus the four constructed schedules. Each candidate is simulated once per season with the same weather, soil, cultivar, spin-up and objective as the PPO environment.

\emph{Without context:} (1) four \emph{constructed reference schedules}: sowing October 15 rainfed, and sowing October 15, September 25, or November 10 with the irrigation rule set to trigger 45 mm, 18 mm per event, cap 120 mm. These dates and parameters were chosen by the authors to span the action window and a moderate irrigation level; they are not agronomic recommendations. (2) An \emph{optimized fixed rule}: the grid candidate with the highest mean objective over the seasons available before the target seasons (Section~\ref{subdesign}).

\emph{With the same pre-season context:} (3) a \emph{contextual nearest-neighbour rule} ($k$NN): for a target season, the three earlier seasons closest in the standardized context of Eq.~(\ref{eq_context}) are identified, and the grid candidate with the highest mean objective over those three seasons is applied. The rule has one parameter ($k = 3$), fixed before the analysis, and receives exactly the information given to PPO. (4) The PPO contextual policy.

\emph{Retrospective ceiling, never a competitor:} (5) a per-season oracle, the grid candidate with the highest objective in each season, selected with the realized weather of that season. The oracle defines the regret and the gap that adaptation could, at most, close within the grid.

\begin{table}[!h]
\centering
\footnotesize
\caption{Comparative design of policies and baselines. Budgets count DSSAT simulation requests.}
\label{tab_design}
\begin{tabular}{p{3.6cm}p{3.6cm}p{2.4cm}p{3.0cm}p{1.4cm}}
\hline
Comparator & Information used & Action space & DSSAT budget & Adaptive \\
\hline
Constructed schedules (4) & none & fixed & 1 run per season & no \\
Optimized fixed rule & seasons before the target block & grid (602 rules) & 602 per season & no \\
Contextual $k$NN rule & context on $t_0$; seasons before the target block & grid (602 rules) & same grid & yes \\
PPO contextual policy & context on $t_0$; training seasons; validation for checkpoint & continuous, 4 dimensions & 16,000--37,000 requests per run & yes \\
Retrospective oracle (ceiling) & realized weather of the season & grid (602 rules) & same grid & retrospective \\
\hline
\end{tabular}
\end{table}

%%============================================================
\subsection{Proximal Policy Optimization for the one-step decision}
\label{subppo}

\emph{Policy and value function.} The policy $\pi_\phi(a \mid s)$ is a diagonal Gaussian over the four pre-squashing action components, with mean $\mu_\phi(s)$ given by a multilayer perceptron with two hidden layers of 32 units and hyperbolic-tangent activation, and with a state-independent log standard deviation vector initialized at $-0.5$ (standard deviation 0.61). Actions are sampled from the Gaussian and clipped to $[-1,1]$ before decoding. A separate network of the same architecture estimates the value $V_\varphi(s)$, the expected objective of the context under the current policy.

\emph{Advantage.} Because each episode has one step and terminates after the simulation, the return of an episode is its objective $r(s,a)$ and the advantage reduces to
\begin{equation}
\hat{A}(s,a) = r(s,a) - V_\varphi(s).
\label{eq_advantage}
\end{equation}
The discount factor ($\gamma = 0.99$) and the generalized advantage estimator ($\lambda = 0.95$) of the implementation have no effect in this formulation. Advantages are normalized within each minibatch.

\emph{Objective.} With $\rho(\phi) = \pi_\phi(a \mid s)/\pi_{\phi_{\text{old}}}(a \mid s)$ the probability ratio between the updated policy and the policy that collected the data, PPO maximizes the clipped surrogate
\begin{equation}
\mathcal{L}(\phi,\varphi) = \mathbb{E}\Big[\min\big(\rho(\phi)\hat{A},\; \mathrm{clip}(\rho(\phi), 1-\epsilon, 1+\epsilon)\,\hat{A}\big)\Big] - c_v\, \mathbb{E}\Big[\big(V_\varphi(s) - r\big)^2\Big] + c_e\, \mathbb{E}\big[\mathcal{H}(\pi_\phi(\cdot \mid s))\big],
\label{eq_ppo}
\end{equation}
with $\epsilon = 0.2$, $c_v = 0.5$ and $c_e = 0$. Each update collects 512 episodes (8 parallel environments $\times$ 64 episodes), performs 10 epochs of minibatch gradient steps with minibatches of 128 episodes, uses Adam with learning rate $3\times 10^{-4}$, and clips the gradient norm at 0.5.

\emph{Training data and regularization.} Each training episode draws one of the training seasons at random. To limit the memorization of season-specific actions, Gaussian noise with standard deviation 0.75 is added to the standardized context of every training episode; validation and test contexts are not perturbed. The noise level and the network size were chosen on the validation seasons of Experiment 1 in one development comparison with seed 42 (noise 0.25 and two layers of 64 units versus noise 0.75 and two layers of 32 units; Section~\ref{subres_e1}); no other hyperparameter was tuned, and the same configuration was used in Experiment 2.

\emph{Validation, checkpoint and stopping.} Every 2,048 training steps the deterministic policy (the Gaussian mean) is evaluated once on each validation season; DSSAT is deterministic, so repeated evaluations of the same season are not drawn. The checkpoint with the highest mean validation objective is retained. In Experiment 1 training runs for at least 20,000 steps and stops after 12 consecutive evaluations without improvement or at 80,000 steps. In Experiment 2, with more seasons per run and a fixed computational budget, training runs for at least 8,192 steps and stops after 6 evaluations without improvement or at 24,576 steps.

\emph{Parallel simulation.} The eight training environments run in separate processes, each with its own DSSAT working directory and its own copy of the DSSAT data folder, because \texttt{DSSATTools} rebuilds that folder at import time and concurrent processes would otherwise overwrite each other's files.

%%============================================================
\subsection{Objective, water constraint, and risk metric}
\label{subobjective}

\emph{Outcome hierarchy.} \emph{Primary results:} yield (kg ha$^{-1}$), seasonal irrigation (mm), and their paired differences between the contextual policy and each comparator in test seasons (Eq.~\ref{eq_paired}). \emph{Systemic interpretation:} the productivity--water relation, the regret with respect to the retrospective oracle, and the behaviour of the policy across contrasting seasons. \emph{Secondary operational result:} the scalar $r$ of Eq.~(\ref{eq_reward}), the optimization function of the search. \emph{Stability outcomes:} DSSAT failures, action saturation, and variation between seeds.

\emph{Optimization function.} The scalar optimized by the search for one season is
\begin{equation}
r = \frac{Y}{Y_{\text{ref}}} - \kappa\, W,
\label{eq_reward}
\end{equation}
with $Y$ the raw DSSAT grain yield (kg ha$^{-1}$), $Y_{\text{ref}} = 4200$ kg ha$^{-1}$ a normalization constant, $W$ the seasonal irrigation (mm), and $\kappa = 0.0015$ mm$^{-1}$. With these values one millimetre of irrigation must raise yield by at least 6.3 kg ha$^{-1}$ to increase $r$. A DSSAT run that fails to complete receives $r = -2$; no failure occurred in either experiment.

\emph{Removal of the statistical correction.} In a previous diagnostic configuration $Y$ was replaced by a ridge regression of the municipal yield on thirteen predictors, including year trend, sowing day and irrigation, fitted on training seasons at five fixed sowing dates. Because the five rows of each season shared one municipal target, the regression carried no information on the response to sowing date or irrigation; its irrigation coefficients were negative ($-0.106$ per mm; rain--irrigation interaction $-0.361$), and within a season it reduced the spread across sowing dates from up to 1,578 kg ha$^{-1}$ of raw yield to 2--48 kg ha$^{-1}$ of corrected yield. The present configuration uses raw DSSAT yield in Eq.~(\ref{eq_reward}); a constant scale or bias correction would not change the ranking of actions within a season.

\emph{Regret and risk.} Regret is the difference between the objective (or yield) of the retrospective oracle and that of the policy in the same season. The risk metrics are the 10th percentile of yield and CVaR$_{10}$ (mean of the yields at or below the 10th percentile), computed over test seasons (and seeds, for PPO).

%%============================================================
\subsection{Experimental design, temporal evaluation, and uncertainty}
\label{subdesign}

\emph{Unit of analysis.} The independent agronomic unit is the season. A PPO episode, a validation evaluation, or a simulated row of the same season is not an independent unit; a training seed is a replication of the optimization procedure and does not increase the agronomic sample size.

\emph{Experiment 1: station record, chronological split.} The INMET seasons are split chronologically (Table~\ref{tab_temporal_registry_castro}): training 2006--2016 (11 seasons), validation 2017--2020 (4), and test 2021--2024 (4). Five seeds (42--46) are trained. The optimized fixed rule and the $k$NN rule are selected on 2006--2020 for the test comparison and on 2006--2016 for the validation comparison. In an earlier version of the workflow the 2021--2024 block had been evaluated once with a policy that received realized-season weather and an objective based on the statistical correction; the changes introduced afterwards (raw yield, the context of Eq.~\ref{eq_context}, the spin-up and weather window, the 4-day inspection, the trigger range, and the regularization of PPO) were decided on the development seasons, and the block was opened once for the configuration reported here. A context-coverage analysis defined after that opening is reported as a post hoc sensitivity.

\begin{table}[!h]
\centering
\scriptsize
\caption{Temporal split of Experiment 1 (INMET A819).}
\label{tab_temporal_registry_castro}
\begin{tabular}{llrrr}
\hline
Block & Years & $n$ & Mean rainfall & Observed yield years \\
      &       &     & (mm season$^{-1}$) & \\
\hline
Training & 2006--2016 & 11 & 892.75 & 11/11 \\
Validation & 2017--2020 & 4 & 1024.75 & 4/4 \\
Test & 2021--2024 & 4 & 980.10 & 4/4 \\
\hline
\end{tabular}
\end{table}

\emph{Experiment 2: gridded record, rolling origin.} The NASA POWER seasons 1984--2024 are evaluated with an expanding-window (rolling-origin) design registered before any run (\texttt{configs/ROLLING\_ORIGIN\_REGISTRY.md}). For a fold whose test block starts in season $s$, PPO trains on 1984 to $s-6$, selects its checkpoint on $s-5$ to $s-1$, and is evaluated on $s$ to $s+3$; the optimized fixed rule and the $k$NN rule are selected on 1984 to $s-1$ (Table~\ref{tab_folds}). Three seeds (42--44) are trained per fold. The registry initially specified test blocks of two seasons (10 folds); before any test-season result of this experiment was produced, the measured simulation rate (about 1.5--5.7 training steps per second) led to an amendment to blocks of four seasons, which keeps the same 20 test seasons and seeds. The run of the first fold with seed 42, whose training and validation seasons were identical in both versions, was retained and its checkpoint evaluated on 2005--2008.

\begin{table}[!h]
\centering
\scriptsize
\caption{Folds of the rolling-origin evaluation (Experiment 2, NASA POWER).}
\label{tab_folds}
\begin{tabular}{lrrrr}
\hline
Fold & PPO training & Checkpoint selection & Test & Fixed and $k$NN rule selection \\
\hline
1 & 1984--1999 (16) & 2000--2004 & 2005--2008 & 1984--2004 \\
2 & 1984--2003 (20) & 2004--2008 & 2009--2012 & 1984--2008 \\
3 & 1984--2007 (24) & 2008--2012 & 2013--2016 & 1984--2012 \\
4 & 1984--2011 (28) & 2012--2016 & 2017--2020 & 1984--2016 \\
5 & 1984--2015 (32) & 2016--2020 & 2021--2024 & 1984--2020 \\
\hline
\end{tabular}
\end{table}

\emph{Paired effects.} For each season $y$, seed $s$, metric $M$, and comparator $b$, the paired effect is
\begin{equation}
\Delta M_{y,s,b} = M_{\pi,y,s} - M_{b,y},
\label{eq_paired}
\end{equation}
where positive values indicate higher yield, objective, or irrigation for PPO.

\emph{Decomposition and intervals.} In Experiment 1 the paired yield differences are decomposed as
\begin{equation}
\Delta_{y,s} = \mu + v_y + w_s + \varepsilon_{y,s},
\label{eq_mixed}
\end{equation}
where $\mu$ is the mean paired difference, $v_y$ the season effect, $w_s$ the seed effect and $\varepsilon_{y,s}$ the residual; with four test seasons, the components are reported as standard deviations and all paired effects are shown, without interval inference. In Experiment 2 the paired effects are first averaged over seeds within each season, and a 95\% interval for the mean over the 20 test seasons is obtained from 5,000 bootstrap resamples of seasons; the between-seed standard deviation of the mean paired difference is reported separately.

\emph{Value decomposition.} The total gain over the agronomic calendar (rainfed October 15) is decomposed as the gain from choosing a good fixed rule (optimized fixed rule minus calendar), the gain from context (contextual $k$NN minus optimized fixed rule), and the incremental value of the algorithm (PPO minus contextual $k$NN). The share of the oracle gain captured by a rule is its gain over the calendar divided by the gain of the oracle over the calendar.

\emph{Skill of the context and seasons with irrigation response.} The association between each context variable and three season outcomes of the candidate grid (best rainfed sowing offset; irrigation gain, defined as the yield of the best grid rule minus that of the best rainfed rule; and yield of the rainfed October 15 schedule) is summarized by Spearman correlations. A season is classed as having an irrigation response when its irrigation gain exceeds 100 kg ha$^{-1}$; this class is used only to describe results and does not enter any policy or comparator.

%%============================================================
\subsection{Reproducibility and configuration}
\label{subrepro}

All code, configurations, raw input files, generated configurations, validation curves, TensorBoard logs, candidate grids and result tables are in the repository \url{https://github.com/GabyQueiroz/RLDSSAT_Soybean}. Experiment 1 is defined by \texttt{configs/experiment\_castro\_t0.yaml} and run by \texttt{run\_final\_castro\_t0.ps1}, which trains the five seeds (\texttt{src/run\_multiseed.py}), collects their evaluations (\texttt{src/summarize\_multiseed.py}) and evaluates the comparators (\texttt{src/final\_comparators.py}); \texttt{src/article\_results.py} produces its tables and figures and \texttt{src/context\_qc\_sensitivity.py} the post hoc context analysis. Experiment 2 is defined by \texttt{configs/experiment\_castro\_power.yaml} and the registry \texttt{configs/ROLLING\_ORIGIN\_REGISTRY.md}; the command \texttt{python -m src.rolling\_origin} with the arguments \texttt{grid}, \texttt{train}, \texttt{reevaluate} and \texttt{analyse} evaluates the candidate grid over 1984--2024, trains the folds (configurations written to \texttt{configs/\_generated\_rolling\_power}), and computes comparators, paired effects and bootstrap intervals; \texttt{src/rolling\_article\_results.py} produces its tables and figures, and \texttt{src/power\_station\_agreement.py} the comparison of the two weather series.

The DSSAT log records DSSAT-CSM Ver. 4.8.0.028; the Python stack includes \texttt{DSSATTools} 3.0.2, Stable-Baselines3 2.7.1, PyTorch 2.10.0, Gymnasium 1.2.2, pandas 2.3.3, NumPy 1.26.4, Matplotlib 3.10.9, and TensorBoard 2.20.0. The runs were executed on a laptop with a 10-core Intel Core i7-1355U processor and eight simulation processes. A training run took 48--56 min in Experiment 1 and 42--109 min in Experiment 2; the candidate grid of Experiment 2 (24,682 simulations) took 27 min with eleven processes.

\begin{figure}[!h]
\centering
\resizebox{\linewidth}{!}{%
\begin{tikzpicture}[
    font=\scriptsize,
    box/.style={rectangle, rounded corners=1pt, draw=black, line width=0.4pt, align=center, text width=2.7cm, minimum height=0.6cm, inner sep=2pt},
    data/.style={box, fill=green!8},
    proc/.style={box, fill=blue!7},
    dec/.style={box, fill=red!7},
    outp/.style={box, fill=gray!10},
    arrow/.style={-Latex, line width=0.45pt},
    darrow/.style={-Latex, line width=0.45pt, dashed}
]

\node[data] (weather) at (0,0) {INMET A819 hourly (Exp.~1) or NASA POWER daily (Exp.~2); NOAA ONI},
\node[proc] (qc) at (0,-1.35) {Quality control and daily series},
\node[data] (field) at (4.4,0) {Synthetic soil profile and generic cultivar 990005},
\node[proc] (calib) at (4.4,-1.35) {DSSAT-CROPGRO configuration (not field-calibrated)},
\node[data] (sidra) at (8.8,0) {SIDRA municipal yield},
\node[outp] (verif) at (8.8,-1.35) {Contextual information only (not an input to DSSAT or to the policy)},
\draw[arrow] (weather) -- (qc),
\draw[arrow] (field) -- (calib),
\draw[darrow] (sidra) -- (verif),
\node[proc] (init) at (0,-3.1) {August 1: soil-water spin-up from drained upper limit},
\node[dec] (t0) at (4.4,-3.1) {Decision date $t_0$ = Sep 1: standardized context $s$ (Eq.~\ref{eq_context})},
\node[dec] (policy) at (8.8,-3.1) {PPO policy or comparator $a=\pi(s)$},
\node[dec] (action) at (13.2,-3.1) {Action $[d,\theta,l,L]$; $l<5$ mm = rainfed},
\node[proc] (dssat) at (13.2,-4.85) {DSSAT season with causal irrigation schedule (4-day inspection)},
\node[outp] (out) at (8.8,-4.85) {Outcomes: raw yield, irrigation, $r$, risk},
\node[outp] (eval) at (4.4,-4.85) {Paired evaluation on test seasons (Exp.~1: 2021--2024; Exp.~2: 2005--2024, rolling origin)},
\draw[arrow] (init) -- (t0),
\draw[arrow] (t0) -- (policy),
\draw[arrow] (policy) -- (action),
\draw[arrow] (action) -- (dssat),
\draw[arrow] (dssat) -- (out),
\draw[arrow] (out) -- (eval),
\draw[darrow] (out.north) -- node[right,font=\tiny]{PPO update (training seasons); checkpoint (validation)} (policy.south),
\draw[arrow] (qc.west) -- ++(-0.35,0) |- ([yshift=-0.75cm]dssat.south) -- (dssat.south),
\draw[arrow] (calib.south) -- ++(0,-0.25) -| ([xshift=0.9cm]dssat.south) ,
\draw[line width=0.6pt] (-1.5,-2.35) -- (14.7,-2.35),
\node[font=\tiny, anchor=south west] at (-1.5,-2.33) {information available before $t_0$},
\node[font=\tiny, anchor=south east] at (14.7,-2.33) {information observed during the season (simulator and causal irrigation rule)},
\end{tikzpicture}}
\caption{System and decision timeline of the framework. Weather from the station (Experiment 1) or from NASA POWER (Experiment 2) passes through quality control to the context and to DSSAT; soil and cultivar are generic and not calibrated with field data; the SIDRA municipal series enters neither DSSAT nor the policy. The decision is taken on $t_0$ from the standardized pre-season context, the season is simulated after a soil-water spin-up from August 1, the irrigation rule is causal in time, PPO is updated on the training seasons and its checkpoint is selected on the validation seasons, and the paired evaluation uses seasons not used for development.}
\label{fig:flowchart_dssat_ppo}
\end{figure}

% =====================================================================
\section{Results}
\label{resultados}

The results are presented in the order of the evidence: the agreement between the two weather series and the response surface of the simulated system, Experiment 1 with the station record, Experiment 2 with the gridded record, and the analysis of the pre-season context common to both.

%%============================================================
\subsection{Weather series and response surface}
\label{subres_landscape}

\emph{Agreement between NASA POWER and the station.} On the days on which the INMET record was valid, daily mean temperature agreed closely between the two series (Pearson $r = 0.96$; station 0.86~$^{\circ}$C cooler than POWER), followed by minimum and maximum temperature ($r = 0.91$ and 0.88), radiation ($r = 0.86$; station 2.5 MJ m$^{-2}$ d$^{-1}$ lower) and daily rainfall ($r = 0.68$; mean absolute difference 3.2 mm d$^{-1}$) (Table~\ref{tab_power_agreement}). For the ten seasons with complete station rainfall, September--April totals differed by 237 mm on average and their correlation was 0.26 (Fig.~\ref{fig_power_inmet}). The gridded series therefore reproduces the temperature regime and the day-to-day rainfall occurrence of the station, but not the seasonal rainfall totals of a single gauge.

\begin{table}[!h]
\centering
\scriptsize
\caption{Agreement between INMET A819 and NASA POWER on days with valid station data. Difference: station minus POWER.}
\label{tab_power_agreement}
\begin{tabular}{lrrrrrr}
\hline
Variable & Days & Station mean & POWER mean & Mean difference & MAE & $r$ \\
\hline
Rainfall (mm d$^{-1}$) & 5,149 & 3.64 & 3.98 & $-0.34$ & 3.22 & 0.68 \\
Maximum temperature ($^{\circ}$C) & 5,532 & 24.24 & 24.60 & $-0.36$ & 1.70 & 0.88 \\
Minimum temperature ($^{\circ}$C) & 5,532 & 12.92 & 13.46 & $-0.54$ & 1.40 & 0.91 \\
Mean temperature ($^{\circ}$C) & 5,532 & 17.49 & 18.35 & $-0.86$ & 1.10 & 0.96 \\
Radiation (MJ m$^{-2}$ d$^{-1}$) & 5,671 & 14.71 & 17.25 & $-2.54$ & 3.12 & 0.86 \\
Sep--Apr rainfall (mm season$^{-1}$) & 10 seasons & 978 & 1,099 & $-121$ & 237 & 0.26 \\
\hline
\end{tabular}
\end{table}

\begin{figure}[!h]
\centering
\includegraphics[width=0.48\linewidth]{fig_power_vs_inmet_season_rain.png}
\caption{September--April rainfall of the INMET A819 station and of NASA POWER for the seasons with at least 90\% valid station rainfall days. Dashed line: 1:1.}
\label{fig_power_inmet}
\end{figure}

\emph{Response surface.} With raw DSSAT yield, the sowing date and irrigation produced large differences within a season. In the fifteen INMET development seasons the range of rainfed yield across the thirteen sowing offsets was 1,853 kg ha$^{-1}$ on average (880--3,932 kg ha$^{-1}$), and irrigation raised the best attainable yield by more than 100 kg ha$^{-1}$ in three seasons (2007, +2,820; 2013, +2,432; 2014, +466 kg ha$^{-1}$) (Fig.~\ref{fig_rainfed_response}). In the 41 NASA POWER seasons, irrigation raised the best attainable yield by more than 100 kg ha$^{-1}$ in 16 seasons, six of them in the 2005--2024 test period (2005, +823; 2008, +388; 2013, +1,373; 2018, +365; 2021, +692; 2023, +1,728 kg ha$^{-1}$). The previous statistical correction had compressed the within-season spread to 2--48 kg ha$^{-1}$, which left the search without a signal to separate sowing dates or irrigation rules.

\begin{figure}[!h]
\centering
\includegraphics[width=0.62\linewidth]{fig_rainfed_sowing_response.png}
\caption{Rainfed DSSAT yield as a function of the sowing offset after September 15 in the INMET development seasons. Thin lines are seasons; thick lines are the means of the training (2006--2016) and validation (2017--2020) seasons.}
\label{fig_rainfed_response}
\end{figure}

%%============================================================
\subsection{Experiment 1: station record}
\label{subres_e1}

\emph{Learning behaviour.} In the five seeds, the mean objective of the training episodes rose from 0.93--0.99 to 1.13--1.16, while the deterministic validation objective reached its maximum within 6,144--12,288 steps and then declined towards the level of the optimized fixed rule; early stopping ended training at 30,720--36,864 steps (Fig.~\ref{fig_learning_curves}). The standard deviation of the policy decreased from 0.61 to 0.44--0.47, the approximate Kullback--Leibler divergence between successive policies remained between 0.004 and 0.027, and the explained variance of the value network rose to 0.31--0.45 (Fig.~\ref{fig_ppo_diagnostics}). In the development comparison with seed 42, the configuration with context noise 0.25 and two layers of 64 units reached 1.187 on the training seasons, close to the oracle (1.213), but 1.164 on validation; the final configuration (noise 0.75, two layers of 32 units) reached 1.228 (Fig.~\ref{fig_regularization}).

\begin{figure}[!h]
\centering
\includegraphics[width=\linewidth]{fig_learning_curves.png}
\caption{Experiment 1: learning curves of PPO for seeds 42--46. (a) Mean objective of training episodes (stochastic actions, perturbed context) after each update of 512 episodes. (b) Objective of the deterministic policy on the four validation seasons, every 2,048 steps; stars mark the selected checkpoints. Horizontal lines: validation objective of the oracle, the rainfed October 15 schedule, and the fixed rule optimized on the training seasons.}
\label{fig_learning_curves}
\end{figure}

\begin{figure}[!h]
\centering
\includegraphics[width=\linewidth]{fig_ppo_diagnostics.png}
\caption{Experiment 1: PPO optimization diagnostics for seeds 42--46: standard deviation of the Gaussian policy, entropy loss, approximate Kullback--Leibler divergence between successive policies, and explained variance of the value network.}
\label{fig_ppo_diagnostics}
\end{figure}

\begin{figure}[!h]
\centering
\includegraphics[width=0.62\linewidth]{fig_regularization_seed42.png}
\caption{Experiment 1: validation objective of seed 42 with the development configuration (context noise 0.25, two hidden layers of 64 units) and with the final configuration (context noise 0.75, two hidden layers of 32 units).}
\label{fig_regularization}
\end{figure}

\emph{Performance by block.} PPO exceeded the optimized fixed rule on the training seasons (1.156 vs 1.130) and on validation (1.226 vs 1.175), but not on the 2021--2024 test block, where it reached 0.961 against 1.134 for the fixed rule selected on 2006--2020 (Table~\ref{tab_by_split}). The paired yield difference with the fixed rule was $-843$ kg ha$^{-1}$ and negative in all four test seasons; the standard deviations of the season, seed and residual components of Eq.~(\ref{eq_mixed}) were 864, 168 and 274 kg ha$^{-1}$. Against the rainfed October 15 schedule and the $k$NN rule, the paired differences were $-279$ and $-118$ kg ha$^{-1}$. The fixed rule captured 67\% of the gain of the oracle over the calendar.

\begin{table}[!h]
\centering
\scriptsize
\caption{Experiment 1: mean objective $r$, yield and irrigation by block. Training and validation comparators are selected on the training seasons; test comparators on 2006--2020. PPO values are means over five seeds (range of seed means in parentheses).}
\label{tab_by_split}
\begin{tabular}{llrrr}
\hline
Block & Policy & $r$ & Yield (kg ha$^{-1}$) & Irrigation (mm) \\
\hline
Training (11) & PPO & 1.156 (1.150--1.162) & 5,150 & 47.0 \\
              & Optimized fixed rule & 1.130 & 5,040 & 46.8 \\
              & Rainfed Oct 15 & 1.020 & 4,284 & 0.0 \\
              & Retrospective oracle & 1.213 & 5,325 & 36.4 \\
Validation (4) & PPO & 1.226 (1.221--1.229) & 5,150 & 0.0 \\
              & Optimized fixed rule & 1.175 & 4,935 & 0.0 \\
              & Contextual $k$NN rule & 1.181 & 4,960 & 0.0 \\
              & Rainfed Oct 15 & 1.221 & 5,127 & 0.0 \\
              & Retrospective oracle & 1.246 & 5,233 & 0.0 \\
Test (4)      & PPO & 0.961 (0.923--0.998) & 4,061 & 3.9 \\
              & Optimized fixed rule & 1.134 & 4,904 & 22.5 \\
              & Contextual $k$NN rule & 0.995 & 4,179 & 0.0 \\
              & Rainfed Oct 15 & 1.033 & 4,340 & 0.0 \\
              & Retrospective oracle & 1.184 & 5,137 & 26.3 \\
\hline
\end{tabular}
\end{table}

\emph{Mechanisms of the test loss.} In 2021 and 2022 all seeds sowed between November 15 and December 3, while the fixed rule sowed on October 9; PPO lost 630 and 413 kg ha$^{-1}$ relative to the fixed rule. The station recorded no valid hourly temperature in August 2021 and August 2022 (0 of 720 hours in each month); after interpolation the 30-day mean temperature took values of 20.1 and 17.7~$^{\circ}$C, against 13.1--15.6~$^{\circ}$C in the development seasons, which correspond to 7.5 and 4.3 standard deviations of the training seasons. In 2024, the only test season with a yield response to irrigation in the station record, the fixed rule applied 75 mm and produced 6,213 kg ha$^{-1}$, against 3,649 kg ha$^{-1}$ for the rainfed calendar; three of the five PPO seeds applied no water and produced 3,670--3,737 kg ha$^{-1}$. The four validation seasons had required no irrigation, and the checkpoints selected on them had not irrigated on validation (Table~\ref{tab_by_split}).

\emph{Post hoc context sensitivity.} When a feature whose window had fewer than 75\% valid days was set to the training mean, which replaced the four station variables of the 2021--2023 test contexts, the late sowings disappeared and the PPO objective rose from 0.961 to 1.008 (4,248 kg ha$^{-1}$); the paired differences became $-92$ kg ha$^{-1}$ against the rainfed calendar and $-656$ kg ha$^{-1}$ against the fixed rule. The 2024 result did not change.

%%============================================================
\subsection{Experiment 2: NASA POWER record, rolling origin}
\label{subres_e2}

\emph{Learning behaviour and checkpoints.} Across the 15 runs, the selected checkpoints were reached at 4,096--24,576 steps, with validation objectives of 1.251--1.374, and training stopped at 16,384--24,576 steps (Table~\ref{tab_checkpoints_e2}, Fig.~\ref{fig_rolling_learning}). The selected checkpoints exceeded, on validation, the fixed rule selected on the training seasons of the fold in folds 1, 4 and 5 (by 0.003--0.036), were at or below it in fold 2 (1.251--1.262 vs 1.261) and around it in fold 3 (1.317--1.339 vs 1.321), and remained below the oracle in all folds (Fig.~\ref{fig_rolling_learning}). The checkpoints selected in fold 2, whose validation block (2004--2008) contained two seasons with an irrigation response, irrigated 54--57 mm on validation; those of the other folds irrigated 3--27 mm. The mean objective of the training episodes rose from 1.11--1.18 to 1.18--1.24, and the standard deviation of the policy decreased from 0.60 to 0.47--0.52 (Fig.~\ref{fig_rolling_diag}). The explained variance of the value network remained between $-0.26$ and 0.20 at the end of training: with 16--32 training seasons and context noise, the critic predicted little of the season-to-season variation of the objective, and the policy updates relied mainly on the minibatch-normalized advantages.

\begin{table}[!h]
\centering
\scriptsize
\caption{Experiment 2: selected checkpoints by fold and seed. Validation values are means over the five validation seasons of the fold.}
\label{tab_checkpoints_e2}
\begin{tabular}{lrrrrrr}
\hline
Test block & Seed & Selected step & Valid. objective & Valid. yield & Valid. irrigation & Stop step \\
           &      &               &                  & (kg ha$^{-1}$) & (mm) & \\
\hline
2005--2008 & 42 & 14,336 & 1.330 & 5,638 & 8.2 & 24,576 \\
           & 43 & 24,576 & 1.330 & 5,607 & 3.2 & 24,576 \\
           & 44 & 6,144 & 1.331 & 5,668 & 12.4 & 18,432 \\
2009--2012 & 42 & 4,096 & 1.251 & 5,611 & 56.6 & 16,384 \\
           & 43 & 6,144 & 1.253 & 5,602 & 53.6 & 18,432 \\
           & 44 & 4,096 & 1.262 & 5,656 & 56.8 & 16,384 \\
2013--2016 & 42 & 6,144 & 1.339 & 5,709 & 13.6 & 18,432 \\
           & 43 & 6,144 & 1.317 & 5,620 & 14.4 & 18,432 \\
           & 44 & 8,192 & 1.323 & 5,653 & 15.2 & 20,480 \\
2017--2020 & 42 & 10,240 & 1.361 & 5,858 & 22.7 & 22,528 \\
           & 43 & 16,384 & 1.366 & 5,875 & 21.6 & 24,576 \\
           & 44 & 6,144 & 1.354 & 5,795 & 17.0 & 18,432 \\
2021--2024 & 42 & 12,288 & 1.374 & 5,902 & 20.6 & 24,576 \\
           & 43 & 16,384 & 1.341 & 5,786 & 24.2 & 24,576 \\
           & 44 & 16,384 & 1.358 & 5,874 & 27.0 & 24,576 \\
\hline
\end{tabular}
\end{table}

\begin{figure}[!h]
\centering
\includegraphics[width=\linewidth]{fig_rolling_learning_curves.png}
\caption{Experiment 2: validation objective of the deterministic policy for each fold and seed, every 2,048 steps. Horizontal lines: validation objective of the oracle, the rainfed October 15 schedule, and the fixed rule optimized on the training seasons of the fold.}
\label{fig_rolling_learning}
\end{figure}

\begin{figure}[!h]
\centering
\includegraphics[width=0.9\linewidth]{fig_rolling_training_diagnostics.png}
\caption{Experiment 2: mean objective of training episodes and standard deviation of the Gaussian policy for the 15 runs, coloured by fold.}
\label{fig_rolling_diag}
\end{figure}

\emph{Test performance.} Over the 20 test seasons, the optimized fixed rule reached a mean objective of 1.308 and a mean yield of 5,683 kg ha$^{-1}$ with 30.3 mm of irrigation; PPO reached 1.290 and 5,590 kg ha$^{-1}$ with 27.4 mm (seed means 1.280--1.307); the $k$NN rule reached 1.275 and 5,506 kg ha$^{-1}$; the rainfed October 15 schedule 1.244 and 5,224 kg ha$^{-1}$; and the oracle 1.376 and 5,878 kg ha$^{-1}$ (Table~\ref{tab_test_e2}). The fixed rule selected in four of the five folds was sowing on October 17 with $\theta = 90$ mm, $l = 15$ mm and $L = 80$ mm; in the fold tested on 2009--2012 the sowing date was October 25. The rainfed calendar had the lowest 10th percentile (3,926 kg ha$^{-1}$) and CVaR$_{10}$ (3,525 kg ha$^{-1}$); the 10th percentile of PPO (5,183 kg ha$^{-1}$) was close to that of the fixed rule (5,238 kg ha$^{-1}$), and its CVaR$_{10}$ (4,639 kg ha$^{-1}$) was lower, mainly because of 2008.

\begin{table}[!h]
\centering
\scriptsize
\caption{Experiment 2: performance over the 20 test seasons (2005--2024). PPO statistics are computed over 60 season--seed pairs; comparator statistics over 20 seasons.}
\label{tab_test_e2}
\begin{tabular}{lrrrrr}
\hline
Policy & $r$ & Yield & P10 & CVaR$_{10}$ & Irrigation \\
       &     & (kg ha$^{-1}$) & (kg ha$^{-1}$) & (kg ha$^{-1}$) & (mm) \\
\hline
Retrospective oracle & 1.376 & 5,878 & 5,539 & 5,353 & 15.5 \\
Optimized fixed rule & 1.308 & 5,683 & 5,238 & 5,092 & 30.3 \\
PPO (3 seeds $\times$ 5 folds) & 1.290 & 5,590 & 5,183 & 4,639 & 27.4 \\
Contextual $k$NN rule & 1.275 & 5,506 & 5,068 & 4,828 & 24.3 \\
Irrigated Oct 15 (45/18/120) & 1.269 & 5,810 & 5,438 & 5,254 & 76.2 \\
Rainfed Oct 15 & 1.244 & 5,224 & 3,926 & 3,525 & 0.0 \\
Irrigated Sep 25 (45/18/120) & 1.238 & 5,697 & 5,374 & 5,141 & 78.6 \\
Irrigated Nov 10 (45/18/120) & 1.196 & 5,511 & 5,259 & 5,032 & 77.4 \\
\hline
\end{tabular}
\end{table}

\emph{Paired effects.} The mean paired difference between PPO and the optimized fixed rule was $-0.018$ in objective (95\% season-bootstrap interval $-0.046$ to 0.008) and $-92$ kg ha$^{-1}$ in yield ($-226$ to 21), with PPO ahead in 10 of the 20 seasons (Table~\ref{tab_paired_e2}, Fig.~\ref{fig_rolling_paired}). Against the $k$NN rule the differences were 0.015 ($-0.021$ to 0.047) and 84 kg ha$^{-1}$ ($-74$ to 232), with PPO ahead in 12 seasons. Against the rainfed October 15 schedule they were 0.046 (0.001 to 0.097) and 366 kg ha$^{-1}$ (106 to 672). PPO had a higher objective than the irrigated September 25 and November 10 schedules (0.052 and 0.094, intervals excluding zero) with about 50 mm less water, and the same objective as the irrigated October 15 schedule (0.021, $-0.022$ to 0.060) with 49 mm less water and 220 kg ha$^{-1}$ less yield. The regret with respect to the oracle was 0.086 ($-288$ kg ha$^{-1}$). The between-seed standard deviation of the mean paired yield difference was 95 kg ha$^{-1}$.

\begin{table}[!h]
\centering
\scriptsize
\caption{Experiment 2: paired effects of PPO over the 20 test seasons (PPO minus comparator; season means over three seeds). Intervals: 95\% bootstrap over seasons (5,000 resamples).}
\label{tab_paired_e2}
\begin{tabular}{lrrrr}
\hline
Comparator & $\Delta r$ (95\% CI) & $\Delta$ yield, kg ha$^{-1}$ (95\% CI) & $\Delta$ irrigation (mm) & Seasons with $\Delta r > 0$ \\
\hline
Optimized fixed rule & $-0.018$ ($-0.046$, 0.008) & $-92$ ($-226$, 21) & $-2.9$ & 10/20 \\
Contextual $k$NN rule & 0.015 ($-0.021$, 0.047) & 84 ($-74$, 232) & 3.1 & 12/20 \\
Rainfed Oct 15 & 0.046 (0.001, 0.097) & 366 (106, 672) & 27.4 & 11/20 \\
Irrigated Oct 15 (45/18/120) & 0.021 ($-0.022$, 0.060) & $-220$ ($-391$, $-83$) & $-48.8$ & 13/20 \\
Irrigated Sep 25 (45/18/120) & 0.052 (0.006, 0.095) & $-106$ ($-278$, 51) & $-51.2$ & 13/20 \\
Irrigated Nov 10 (45/18/120) & 0.094 (0.054, 0.131) & 80 ($-95$, 226) & $-50.0$ & 18/20 \\
Retrospective oracle & $-0.086$ ($-0.122$, $-0.055$) & $-288$ ($-423$, $-164$) & 11.9 & 2/20 \\
\hline
\end{tabular}
\end{table}

\begin{figure}[!h]
\centering
\includegraphics[width=\linewidth]{fig_rolling_paired_yield.png}
\caption{Experiment 2: paired yield differences between PPO (mean over three seeds) and each comparator in the 20 test seasons.}
\label{fig_rolling_paired}
\end{figure}

\emph{Irrigation decisions.} PPO concentrated irrigation in the seasons with an irrigation response (Fig.~\ref{fig_rolling_irrigation}). In the six test seasons with a response it applied 75.2 mm on average and in the fourteen others 6.8 mm, against 78.3 and 9.6 mm for the fixed rule and 62.5 and 7.9 mm for the $k$NN rule; the oracle applied 51.7 mm and none. In the six response seasons PPO produced 1,251 kg ha$^{-1}$ more than the rainfed calendar and 398 kg ha$^{-1}$ less than the fixed rule (objective difference $-0.090$); in the fourteen others it produced 13 kg ha$^{-1}$ less than the calendar and 39 kg ha$^{-1}$ more than the fixed rule (objective difference 0.013). The largest loss relative to the fixed rule occurred in 2008 ($-1,089$ kg ha$^{-1}$), when PPO applied 45 mm and the fixed rule 75 mm. PPO sowed between October 5 and November 6 in all test seasons. By fold, PPO was below the fixed rule in the 2005--2008 block (1.207 vs 1.278), above it in 2009--2012 (1.346 vs 1.319), and within 0.025 of it in the other three blocks (Table~\ref{tab_folds_e2}).

\begin{figure}[!h]
\centering
\includegraphics[width=\linewidth]{fig_rolling_irrigation_by_season.png}
\caption{Experiment 2: seasonal irrigation of PPO (mean over three seeds), the optimized fixed rule, the $k$NN rule and the oracle in each test season. Values in parentheses: irrigation gain of the best grid rule over the best rainfed rule (kg ha$^{-1}$).}
\label{fig_rolling_irrigation}
\end{figure}

\begin{table}[!h]
\centering
\scriptsize
\caption{Experiment 2: mean objective $r$ and yield (kg ha$^{-1}$, in parentheses) by test block.}
\label{tab_folds_e2}
\begin{tabular}{lrrrrr}
\hline
Test block & PPO & Optimized fixed rule & Contextual $k$NN & Rainfed Oct 15 & Oracle \\
\hline
2005--2008 & 1.207 (5,257) & 1.278 (5,614) & 1.247 (5,427) & 1.163 (4,885) & 1.310 (5,626) \\
2009--2012 & 1.346 (5,662) & 1.319 (5,562) & 1.292 (5,521) & 1.373 (5,766) & 1.413 (5,934) \\
2013--2016 & 1.323 (5,703) & 1.347 (5,784) & 1.306 (5,602) & 1.250 (5,252) & 1.392 (5,973) \\
2017--2020 & 1.307 (5,710) & 1.308 (5,778) & 1.287 (5,459) & 1.315 (5,522) & 1.397 (5,866) \\
2021--2024 & 1.267 (5,621) & 1.286 (5,676) & 1.241 (5,520) & 1.118 (4,695) & 1.370 (5,990) \\
\hline
\end{tabular}
\end{table}

\emph{Value decomposition and productivity--water relation.} Relative to the rainfed calendar, the optimized fixed rule added 0.064 to the objective and 459 kg ha$^{-1}$ to yield and captured 48\% of the gain of the oracle over the calendar; PPO captured 35\%. Adding context through the $k$NN rule changed the objective by $-0.033$ relative to the fixed rule, and replacing the $k$NN rule by PPO changed it by 0.015. In the yield--irrigation plane (Fig.~\ref{fig_rolling_yw}), the three PPO seeds lay at 24--34 mm and 5,529--5,700 kg ha$^{-1}$, close to the fixed rule and to the upper envelope of the grid rules, whereas the constructed irrigated schedules used about 76--79 mm.

\begin{figure}[!h]
\centering
\includegraphics[width=0.68\linewidth]{fig_rolling_yield_water.png}
\caption{Experiment 2: mean yield and mean seasonal irrigation over the 20 test seasons for the 598 grid rules (grey), the comparators and the three PPO seeds. The grid points describe the attainable set and are not available to any operational comparator.}
\label{fig_rolling_yw}
\end{figure}

%%============================================================
\subsection{Pre-season context: skill and coverage}
\label{subres_context}

The association between the context variables and the season outcomes was weak in both series. In the fifteen INMET development seasons, the largest Spearman correlation was between 90-day pre-season rainfall and the best rainfed sowing offset ($\rho = -0.61$); no variable was associated with the irrigation gain ($|\rho| \leq 0.24$), and the ONI terms had $|\rho| \leq 0.22$. In the 41 NASA POWER seasons all correlations were $|\rho| \leq 0.33$; the largest was between 90-day radiation and the yield of the rainfed October 15 schedule (0.33), and no variable reached $|\rho| > 0.20$ with the irrigation gain. In the station record, the two test contexts built from interpolated temperature lay 6.45 and 4.09 standardized units from the nearest training season, against 1.38--3.55 among the training seasons, and the radiation of the 2017--2019 validation seasons was 4.2--5.2 standard deviations below the training mean, consistent with sensor degradation. The gridded series has no missing values, and its contexts remained within the range of the earlier seasons.

%%============================================================
\subsection{Run integrity and scope of inference}
\label{subres_generalization}

No DSSAT run failed in the 9,030 development and 2,408 test grid simulations of Experiment 1, in the 24,682 grid simulations of Experiment 2, or in the PPO evaluations. Two sowings selected by PPO in Experiment 1 (December 2 and 3) exceed by one and two days a conservative 150-day horizon ending on April 30; all sowings of Experiment 2 fell between October 5 and November 6. The conclusions are restricted to Castro, to the synthetic profile and the generic cultivar, and to the irrigation-rule family described in Section~\ref{subdecision}.

% =====================================================================
\section{Discussion}
\label{discussao}

The two experiments answer the central question in different ways, and the difference identifies the conditions under which a context-conditioned policy can be evaluated. With the station record (19 seasons, gaps in the pre-season windows, four wet validation seasons), PPO exceeded the fixed rule inside the development seasons and fell below every comparator on four later seasons. With the gridded record (41 seasons without gaps), PPO matched the optimized fixed rule within the season-to-season uncertainty over 20 test seasons, exceeded the rainfed calendar by 366 kg ha$^{-1}$, and used less water than the constructed irrigated schedules for a higher or equal objective. In neither experiment did PPO exceed the optimized fixed rule.

%%============================================================
\subsection{Mechanisms: climate--soil--sowing--water interaction}
\label{subdisc_mechanisms}

Two processes determine the value of an action in the simulated system. The sowing date places the reproductive period in a window of radiation, temperature and rainfall, and its effect is present in every season. Irrigation has value only in the seasons in which the 198 mm of available water of the profile is depleted during the reproductive period: three of fifteen INMET development seasons and sixteen of 41 NASA POWER seasons. A fixed rule with a high trigger exploits this structure without knowing in advance which season will be dry, because the trigger reacts to the in-season water balance. In Experiment 2 the fixed rule irrigated 78 mm on average in the response seasons and 10 mm in the others. PPO reproduced this pattern (75 and 7 mm), so most of its irrigation behaviour was mediated by the same causal trigger rather than by the pre-season context. This is consistent with the observation of \citet{kelly2024assessing} that optimized heuristic irrigation policies were difficult to outperform outside strongly water-limited conditions.

%%============================================================
\subsection{Value of context and information}
\label{subdisc_information}

In both experiments the decomposition attributes the gain over the calendar to the choice of a good fixed rule and no positive contribution to context. The pre-season variables carried little information on the outcome that differed most between seasons: no context variable was associated with the irrigation gain in either series, and the ONI terms, the only seasonal-scale information in the context, were not associated with sowing response or irrigation gain. This agrees with the view that the value of pre-season climate information must be measured for the decision at hand rather than assumed \citep{hansen2002climate}. Where the context has low skill, a policy that conditions on it can at best reproduce the fixed rule, and the $k$NN rule and PPO, which received the same information, ended within 0.015 of each other in Experiment 2.

The record used to train and evaluate the policy changed the result more than the algorithm. With eleven training seasons and gaps in the station record, the policy associated contexts with season-specific actions, extrapolated when interpolated temperature placed two test contexts outside the training range, and inherited from four wet validation seasons a preference for not irrigating. With 16--32 training seasons and a homogeneous gridded series, the same algorithm and hyperparameters produced policies whose test performance stayed within 0.025 of the fixed rule in four of five blocks. The gridded record has its own limitation: its seasonal rainfall agreed poorly with the station gauge, so the 41 seasons describe the climate of a grid cell rather than that of the field.

%%============================================================
\subsection{When an adaptive policy is worth the complexity}
\label{subdisc_value}

The results indicate three conditions for a context-conditioned policy to add value over a fixed rule in a system of this type. The context must have measurable skill for the component of the decision that varies between seasons, here the occurrence of a reproductive-period drought; the training and validation seasons must include enough seasons of each hydrological type for the selection criterion to reward the right behaviour; and the pre-season inputs must be complete or carry a coverage check before they reach the policy. The second and third conditions were met in Experiment 2 and the first was not, and the policy then matched but did not exceed the fixed rule. Where the first condition is not met, a fixed rule whose irrigation component reacts to the season, selected on all available seasons, is the reference against which reinforcement learning for crop management should be reported \citep{balderas2025comparative, noa2026hybrid}. Context variables with demonstrated skill for seasonal water deficit, such as soil-moisture states or calibrated seasonal forecasts, would be the next information to test with the same design.

%%============================================================
\subsection{Resource allocation and risk}
\label{subdisc_resources}

In Experiment 2 the optimized fixed rule and PPO used 27--30 mm of irrigation per season on average and raised mean yield by 366--459 kg ha$^{-1}$ over the rainfed calendar, almost entirely through the response seasons, where rainfed yield fell to 4,086 kg ha$^{-1}$ on average. The constructed schedules used more than twice that water for at most 220 kg ha$^{-1}$ more yield than PPO. The low-yield tail was governed by the response seasons for the rules that did not irrigate and, for PPO, by 2008, when it applied less water than the fixed rule. These are properties of the simulated system with a generic cultivar and a synthetic soil, and they are not an operational recommendation of irrigation depth.

%%============================================================
\subsection{Scope and limitations}
\label{subdisc_limitations}

The scope of inference is Castro only. The main limitations are the synthetic soil and generic cultivar without field calibration, the use of municipal yield only as contextual information, the quality of the station record in recent years, the spatial resolution of the gridded record and its weak agreement with the station on seasonal rainfall, the proxy water balance of the irrigation rule, and the simplified one-step representation of management. In Experiment 1 the test block had been evaluated once in an earlier configuration with a different context and objective. In Experiment 2 the fold length was amended before any test result was produced, the training budget was limited to 24,576 steps per run, and the value network explained little of the variance of the objective. Logistic constraints, irrigation costs beyond the linear penalty, nutrient limitations, cultivar choice, and machinery availability were not modelled.

% =====================================================================
\section{Conclusion}
\label{conclusao}

The study asked whether a context-conditioned sowing-and-irrigation policy adds value over a fixed rule in Castro. With the INMET station record, PPO exceeded the optimized fixed rule on the training and validation seasons but produced 843 kg ha$^{-1}$ less on the 2021--2024 test block, because missing station temperature moved two test contexts far outside the training range and the validation seasons contained no dry season. With the 41-season NASA POWER record and a registered rolling-origin evaluation over 20 test seasons, PPO produced 5,590 kg ha$^{-1}$ with 27.4 mm of irrigation; its paired difference with the optimized fixed rule was $-92$ kg ha$^{-1}$ (95\% interval $-226$ to 21), and it exceeded the rainfed October 15 calendar by 366 kg ha$^{-1}$ (106 to 672) by irrigating in the seasons with a yield response to water.

Regarding the three directional expectations: the value of adaptation depended on water availability, because the differences between rules were concentrated in the seasons in which the soil water was depleted; the effect of soil storage appeared through the high trigger selected for the fixed rule and learned by PPO, which irrigated only when the 198 mm profile was depleted; and the value of pre-season information was low, in line with its weak association with sowing response and irrigation gain in both weather series. The contribution is a reproducible DSSAT--PPO workflow with a comparator structure in which every rule is evaluated on the same seasons, and the evidence that the length, completeness and hydrological composition of the weather record determine whether a learned policy reaches the level of a well-chosen fixed rule. Because DSSAT was not calibrated with local field data, the results should not be interpreted as field-scale recommendations for sowing date or supplemental irrigation.
%===============================================================
\section*{Declaration of generative AI and AI-assisted technologies in the writing process}

The authors used ChatGPT only to support English-language editing and readability checks. The authors reviewed and verified the manuscript content, numerical results, cited literature, code outputs, and interpretations, and take full responsibility for the submitted version.

\section*{Declaration of competing interest}
The authors declare that they have no known financial or personal competing interests that could have influenced the work reported in this article.

\section*{Compliance with ethical standards}
This study did not involve human participants, animals, or sensitive personal data. Therefore, no ethical approval was required.

\section*{Data and code availability}

Code, data, figures, and accompanying materials for the reported computational workflow are available on Zenodo as software, version v1.0.0, at DOI \url{https://doi.org/10.5281/zenodo.22260538}. The repository is available at \url{https://github.com/GabyQueiroz/RLDSSAT_Soybean} and includes the executable computational environment, redistributable data, seeds, logs, scripts used to generate figures and tables, baseline configurations, and license metadata.

\section*{Funding}
This research was funded by national funds through FCT – Fundação para a Ciência e a Tecnologia, I.P., under project UIDB/00048/2020 (DOI: 10.54499/UIDB/00048/2020), by the National Council for Scientific and Technological Development (CNPq, Brazil) through Productivity in Research Scholarships (Grant Nos. 301644/2022-5 and 303909/2025-0), and by the Fundação Araucária through a postdoctoral fellowship (Grant No. 56/2026 FAUEPG).

\section*{CRediT authorship contribution statement}

Marcella Scoczynski Ribeiro Martins: Conceptualization, Investigation, Writing -- original draft, Writing -- review and editing. Gabrielly de Queiroz Pereira: Data curation, Formal analysis, Software, Visualization, Writing -- original draft, Writing -- review and editing. Hilkija Gaius Tosso: Methodology, Validation, Writing -- review and editing. Maria Salete Marcon Gomes Vaz: Supervision, Methodology, Writing -- review and editing. Jose Carlos Ferreira da Rocha: Supervision, Methodology, Writing -- review and editing. Viviane M. Gomes Pacheco: Methodology, Writing -- review and editing. Cloves Goncalves Rodrigues: Methodology, Writing -- review and editing. Antonio Paulo Coimbra: Supervision, Writing -- review and editing. Wesley Pacheco Calixto: Conceptualization, Methodology, Supervision, Project administration, Funding acquisition, Writing -- review and editing.

\section*{Supplementary material}

The supplementary material contains the source tables, diagnostic outputs, and additional analyses that support the results reported in the main text. The main article retains the system and decision-timeline figures, the DSSAT configuration and validation limits, the paired effects against the final comparators, the productivity--water analysis, the temporal inference scope for Castro, and the context--action relationship.

\begin{itemize}
\item \textbf{S1. Configurations.} \texttt{experiment\_castro\_t0.yaml}, \texttt{experiment\_castro\_power.yaml}, the rolling-origin registry with its amendment, per-run generated configurations, and the context scaler of each run.
\item \textbf{S2. Soil profile.} Soil profile used by the executable DSSAT workflow, reported in \texttt{synthetic\_soil\_profile\_s2.csv}.
\item \textbf{S3. Weather data.} INMET completeness and imputation diagnostics, the unmodified NASA POWER file with its API request, the POWER--INMET agreement tables, the ONI file, the SIDRA pairing audit, and planting-window checks.
\item \textbf{S4. Previous diagnostic configuration.} Statistical yield correction, its ablation, and the policy evaluated with realized-season context, retained to document the redesign.
\item \textbf{S5. Learning curves.} Validation curves, TensorBoard scalars, selected checkpoints and stopping steps of every run of Experiments 1 and 2, and the development runs of seed 42.
\item \textbf{S6. Candidate grids.} Simulated yield, irrigation and objective of the 602 grid rules for the INMET seasons 2006--2024 and the NASA POWER seasons 1984--2024.
\item \textbf{S7. Actions and paired effects.} Actions by season and seed, paired effects against all comparators, bootstrap intervals, value decomposition, regret, and water-penalty sensitivity.
\item \textbf{S8. Context analysis.} Context variables by season, Spearman correlations with season outcomes, position of each context relative to the training seasons, and the post hoc coverage sensitivity of Experiment 1.
\end{itemize}


\bibliographystyle{elsarticle-harv}
\bibliography{01_marcella_v3.bib}


% =====================================================================



\end{document}
